\chapter{Internationalisation and locale}
\label{ch:i18n}

\begin{aside}
A TUI library used outside one language and one terminal soon
encounters every weird edge of internationalisation: non-ASCII
text, right-to-left scripts, locale-dependent number
formatting, time zones. \maya{} does not try to solve every i18n
problem (no library does), but it gives you primitives that
play well with the standard C++ approaches. This chapter is the
guide.
\end{aside}

\section{Three i18n axes}

International apps face three orthogonal concerns:

\begin{itemize}[leftmargin=*]
  \item \textbf{Translation.} The text shown to the user is in
    their language.
  \item \textbf{Locale formatting.} Numbers, dates, currencies
    follow local conventions.
  \item \textbf{Text rendering.} Non-Latin scripts, RTL, complex
    glyphs render correctly.
\end{itemize}

Each is a separate problem with its own tools.

\section{Translation}

The canonical pattern: wrap user-visible strings in a function
that looks up the translation:

\begin{lstlisting}
std::string_view _(const char* key);  // looks up by key

auto label = text(_(\"open_file\"));
\end{lstlisting}

\code{\_} (underscore) is the conventional name. The key is a
short identifier; the actual text comes from a translation
table indexed by current locale.

\subsection*{Loading translations}

\begin{lstlisting}
struct Translations {
    std::unordered_map<std::string, std::string> strings;

    static Expected<Translations, std::string>
    load(std::string_view lang_tag);
};
\end{lstlisting}

Translation files are typically JSON or PO format. Load once at
startup; hold in a global (or in the model).

\subsection*{Plurals}

English has two plural forms (one, many); some languages have
six. The standard solution is the CLDR plural rules:

\begin{lstlisting}
std::string_view ngettext(const char* key, int n);
\end{lstlisting}

The function picks the right form based on $n$ and the current
locale's plural rule. \maya{} does not ship one; use a library
(\code{boost::locale}, ICU) or roll your own.

\section{Locale formatting}

Numbers, dates, currencies follow locale conventions. C++26's
\code{std::format} respects the current locale by default:

\begin{lstlisting}
std::string s = std::format(\"{:L}\", 1234567);
// en-US: \"1,234,567\"
// de-DE: \"1.234.567\"
// fr-FR: \"1 234 567\"
\end{lstlisting}

The \code{L} format specifier enables locale awareness.

\subsection*{Setting the locale}

\begin{lstlisting}
std::locale::global(std::locale{\"de_DE.UTF-8\"});
\end{lstlisting}

Setting the global locale affects \code{format}, \code{ostream},
and the C functions that read the locale. Set it once at
startup based on the user's environment.

\subsection*{Time zones}

C++20 introduced \code{std::chrono::zoned\_time}:

\begin{lstlisting}
auto now = std::chrono::system_clock::now();
auto local = std::chrono::zoned_time{std::chrono::current_zone(), now};
std::cout << local << \"\\n\";
\end{lstlisting}

The output respects the system time zone. \maya{} apps that
display timestamps should always use \code{zoned\_time}
through the user's local zone, not the system clock directly.

\section{Text rendering}

\maya{}'s grapheme walker (Chapter~\ref{ch:text}) handles
non-Latin scripts correctly. The width table covers all
codepoints. But two complications remain.

\subsection*{Bidirectional text}

Arabic and Hebrew are read right-to-left. Modern terminals
(GNOME Terminal, Konsole, Alacritty) implement BiDi rendering;
\maya{} relies on this and emits text in logical order.

The catch: \emph{mixed} LTR/RTL inside a single text node is
where edge cases live. If the line contains both English and
Arabic, the visual order depends on the terminal's BiDi
implementation. \maya{} does not attempt to predict it.

\subsection*{Complex shaping}

Devanagari, Arabic ligatures, vowel marks --- these require
the font shaper to combine sequences into glyphs. The terminal's
font handles this, not \maya{}. The grapheme walker counts the
cells correctly; the terminal renders the glyphs.

\section{Identifying the locale}

\begin{lstlisting}
std::string detect_locale() {
    if (auto* l = std::getenv(\"LANG\")) return l;
    if (auto* l = std::getenv(\"LC_ALL\")) return l;
    return \"C\";
}
\end{lstlisting}

POSIX uses \code{LANG} and \code{LC\_*} variables. Win32 uses
\code{GetUserDefaultLocaleName}. The platform abstraction can
unify them.

\section{Number input}

When parsing user-typed numbers, the locale matters. \code{1,5}
is one and a half in German, fifteen in English. The cleanest
approach: pick a canonical form for storage, convert only at
display.

\begin{lstlisting}
double parse_locale_number(std::string_view s);  // tolerant
std::string format_locale_number(double n);       // locale-aware

// Internally: always double.
double value = parse_locale_number(user_input);
display = format_locale_number(value);
\end{lstlisting}

The model holds the unambiguous form; the view formats for
display.

\section{Plural-aware UIs}

For status messages with counts:

\begin{lstlisting}
auto msg = std::format(\"{} {}\",
                       std::format(\"{:L}\", count),
                       ngettext(\"file_singular\", \"file_plural\", count));
\end{lstlisting}

The translation table has separate entries for the singular and
plural forms. The format pulls the right one.

\section{The C locale}

The default C locale (\code{LC\_ALL=C}) is the lowest common
denominator: ASCII only, no comma separators, English. For
servers and scripts, this is intentional --- predictable.

For an interactive TUI, default to the user's locale. Honour
\code{LC\_ALL=C} as an explicit opt-out.

\section{Right-to-left UI shells}

If the entire UI is in Arabic or Hebrew, you may want the
\emph{layout} flipped (sidebar on the right). This is
deliberate; the flexbox model supports it via
\code{Direction::RowReverse}:

\begin{lstlisting}
auto root = h(...) | (rtl ? direction<RowReverse> : direction<Row>);
\end{lstlisting}

Each row places its children in reverse order; columns are
unaffected.

\section{Pitfalls}

\begin{warning}
\textbf{Hard-coded English in widget defaults.} If your
widget shows ``OK'' and ``Cancel'', those strings need
translation. Avoid hard-coded UX strings; route them through
\code{\_(\ldots)}.
\end{warning}

\begin{warning}
\textbf{Assuming character width.} A name field that allocates
20 characters works for English ``Alice'' but fails for the
six-cell ``\\verb|ya-ma-da|'' (a CJK transliteration). Allocate
by cell, not by character count.
\end{warning}

\begin{warning}
\textbf{Date formatting in strings.} Hard-coding ``MM/DD/YYYY''
breaks for users who expect ``DD.MM.YYYY''. Use
\code{std::format} with a locale-aware specifier.
\end{warning}

\begin{warning}
\textbf{Truncating at byte boundaries.} A UTF-8 truncation at
byte $n$ might split a multi-byte character. Always truncate at
grapheme boundaries (Chapter~\ref{ch:text}).
\end{warning}

\section*{Workshop}

\begin{enumerate}[leftmargin=*]
  \item Add a small translation table to a counter app. Switch
    locales at runtime; verify the strings change.
  \item Format a number with \code{std::format}'s \code{L}
    specifier in three locales. Note the differences.
  \item Type Arabic text in a text input widget. Verify the
    cells are counted correctly even if the visual order is
    terminal-dependent.
\end{enumerate}

\begin{recap}
i18n is three problems: translation, locale formatting, text
rendering. \code{std::format} handles locale numbers and dates;
the grapheme walker handles non-Latin text; the terminal handles
shaping and BiDi. \maya{} stays out of the way; you compose the
right tools. Default to the user's locale, opt out only when you
need predictability.
\end{recap}
